\documentclass[11pt,a4paper]{article}

% ─── Pacchetti ────────────────────────────────────────────────────────────────
\usepackage[T1]{fontenc}
\usepackage[utf8]{inputenc}
\usepackage[italian]{babel}
\usepackage{microtype}
\usepackage{geometry}
\geometry{a4paper, top=25mm, bottom=25mm, left=25mm, right=25mm}

\usepackage{amsmath, amssymb, amsthm}
\usepackage{bm}          % bold math
\usepackage{physics}     % \dv, \pdv, \norm, ecc.
\usepackage{siunitx}
\sisetup{
  per-mode        = symbol,
  detect-all,
  exponent-product = \cdot,
  output-decimal-marker = {.}
}

\usepackage{booktabs}
\usepackage{array}
\usepackage{multirow}
\usepackage{xcolor}
\usepackage{graphicx}
\usepackage{float}
\usepackage{caption}
\usepackage{subcaption}
\usepackage{hyperref}
\hypersetup{
  colorlinks = true,
  linkcolor  = blue!60!black,
  citecolor  = green!40!black,
  urlcolor   = blue!70!black
}

% TikZ
\usepackage{tikz}
\usepackage{pgfplots}
\pgfplotsset{compat=1.18}
\usetikzlibrary{
  arrows.meta,
  decorations.pathreplacing,
  patterns,
  positioning,
  calc,
  shapes.geometric,
  fit
}

\usepackage{listings}
\lstset{
  basicstyle   = \ttfamily\small,
  keywordstyle = \color{blue!70!black}\bfseries,
  commentstyle = \color{green!50!black}\itshape,
  stringstyle  = \color{red!60!black},
  breaklines   = true,
  frame        = single,
  numbers      = left,
  numberstyle  = \tiny\color{gray}
}

% ─── Nuovi comandi ────────────────────────────────────────────────────────────
\newcommand{\tens}[1]{\bm{#1}}
\newcommand{\epdot}{\dot{\varepsilon}^{p}}
\newcommand{\sdot}{\dot{s}}
\newcommand{\sstar}{s^{*}}
\newcommand{\sac}{SAC305}
\newcommand{\comsol}{\textsc{Comsol}\,Multiphysics}

% ─── Intestazione ─────────────────────────────────────────────────────────────
\title{%
  \textbf{Modello Viscoplastico di Anand per lo Strato di Solder \sac{}}\\[4pt]
  \large Teoria, Parametri e Implementazione in \comsol{}\\
  \large per Moduli di Potenza SiC a Media Tensione (SSCB 1--10\,kV)
}
\author{Davide Bagnara\\
  \small Power Electronics \& Control Systems
}
\date{\today}

% ══════════════════════════════════════════════════════════════════════════════
\begin{document}
\maketitle
\tableofcontents
\bigskip
\hrule
\bigskip

% ══════════════════════════════════════════════════════════════════════════════
\section{Introduzione e Contesto}
% ══════════════════════════════════════════════════════════════════════════════

Nei moduli di potenza a base SiC destinati ad applicazioni SSCB (\emph{Solid-State
Circuit Breaker}) a media tensione (1--10\,kV, $<$\,1\,kA), lo strato di solder
\sac{} (Sn\,96.5 / Ag\,3 / Cu\,0.5) costituisce l'interfaccia meccanica e termica
tra il die semiconduttore e il substrato DBC (\emph{Direct Bonded Copper}).

Durante un evento di fault, la corrente raggiunge valori di picco di 2--3$\times
I_{\mathrm{nom}}$ in un tempo $t_{\mathrm{open}} = \SI{2}{\micro\second}$--\SI{5}{\micro\second},
generando impulsi termici rapidi e ciclici nel die SiC. Il differenziale di CTE
tra il die ($\alpha_{\mathrm{SiC}} \approx 4\;\mu\mathrm{m/(m\cdot K)}$) e il solder
($\alpha_{\mathrm{SAC}} \approx 21\;\mu\mathrm{m/(m\cdot K)}$) produce deformazioni
inelastiche accumulate che determinano la vita a fatica del modulo.

Il modello costitutivo di \textbf{Anand} \cite{Anand1985} è quello più utilizzato
per descrivere il comportamento viscoplastico del solder a temperatura: esso
incorpora creep e plasticità in un'unica formulazione e non richiede la separazione
dei meccanismi di deformazione.

% ══════════════════════════════════════════════════════════════════════════════
\section{Decomposizione della Deformazione}
% ══════════════════════════════════════════════════════════════════════════════

La deformazione totale $\tens{\varepsilon}_{\mathrm{tot}}$ si decompone in tre
contributi additivi (ipotesi di piccole deformazioni):
%
\begin{equation}
  \tens{\varepsilon}_{\mathrm{tot}}
  = \tens{\varepsilon}^{\mathrm{el}}
  + \tens{\varepsilon}^{\mathrm{th}}
  + \tens{\varepsilon}^{vp}
  \label{eq:decomp}
\end{equation}
%
con:
\begin{align}
  \tens{\varepsilon}^{\mathrm{el}} &= \mathbf{C}^{-1} : \tens{\sigma}
  \quad \text{(elastica, legge di Hooke)} \label{eq:el}\\
  \tens{\varepsilon}^{\mathrm{th}} &= \alpha_{\mathrm{CTE}}(T-T_{\mathrm{ref}})\,\mathbf{I}
  \quad \text{(termica, isotropa)} \label{eq:th}\\
  \tens{\varepsilon}^{vp} &= \int_0^t \dot{\tens{\varepsilon}}^{vp}\,\dd t'
  \quad \text{(viscoplastica, Anand)} \label{eq:vp}
\end{align}

Il tensore di rigidezza $\mathbf{C}$ dipende da $E(T)$ e $\nu$
(v.~Sezione~\ref{sec:props}).

% ══════════════════════════════════════════════════════════════════════════════
\section{Equazioni Costitutive del Modello di Anand}
% ══════════════════════════════════════════════════════════════════════════════

\subsection{Variabili di stato}

Il modello di Anand è \textbf{unificato}: non distingue plasticità istantanea
da creep. Le variabili di stato sono:
\begin{itemize}
  \item $\sigma$ -- tensione equivalente di von Mises [\si{\pascal}]
  \item $s$ -- \emph{deformation resistance}, variabile interna scalare [\si{\pascal}]
  \item $T$ -- temperatura assoluta [\si{\kelvin}]
  \item $\epdot$ -- velocità di deformazione plastica equivalente [\si{\per\second}]
\end{itemize}

\subsection{Legge di flusso plastico}

\begin{equation}
  \boxed{
    \epdot = A\,\exp\!\left(-\frac{Q}{RT}\right)
    \left[\sinh\!\left(\xi\,\frac{\sigma}{s}\right)\right]^{1/m}
  }
  \label{eq:flow}
\end{equation}

\begin{table}[H]
  \centering
  \caption{Parametri della legge di flusso (Eq.~\ref{eq:flow})}
  \begin{tabular}{clll}
    \toprule
    Simbolo & Significato & Unità & Nota \\
    \midrule
    $A$     & Costante pre-esponenziale & \si{\per\second} & fattore di frequenza \\
    $Q$     & Energia di attivazione & \si{\joule\per\mol} & termine Arrhenius \\
    $R$     & Costante dei gas & \SI{8.314}{\joule\per\mol\per\kelvin} & costante \\
    $\xi$   & Stress multiplier & -- & sensitività alla tensione \\
    $m$     & Strain rate sensitivity & -- & $0 < m \leq 1$ \\
    $s$     & Resistenza corrente & \si{\pascal} & evolve con Eq.~\ref{eq:sdot} \\
    \bottomrule
  \end{tabular}
  \label{tab:flow_params}
\end{table}

Il termine $\sinh(\xi\,\sigma/s)$ descrive tre regimi:
\begin{align}
  \xi\,\frac{\sigma}{s} \ll 1 &\implies \epdot \propto \left(\frac{\sigma}{s}\right)^{1/m}
  \quad \text{(regime Norton)} \label{eq:norton}\\
  \xi\,\frac{\sigma}{s} \gg 1 &\implies \epdot \propto \exp\!\left(\frac{\xi\sigma}{ms}\right)
  \quad \text{(regime esponenziale)} \label{eq:exp}
\end{align}

\subsection{Evoluzione della variabile interna \texorpdfstring{$s$}{s}}

\begin{equation}
  \boxed{
    \sdot = h_0\,\left|1 - \frac{s}{\sstar}\right|^{a}
    \operatorname{sign}\!\left(1 - \frac{s}{\sstar}\right)\cdot\epdot
  }
  \label{eq:sdot}
\end{equation}

con la resistenza di saturazione dipendente da $T$ e $\epdot$:
\begin{equation}
  \boxed{
    \sstar = \hat{s}\left[\frac{\epdot}{A}\,\exp\!\left(\frac{Q}{RT}\right)\right]^{n}
  }
  \label{eq:sstar}
\end{equation}

\begin{table}[H]
  \centering
  \caption{Parametri dell'equazione di evoluzione (Eq.~\ref{eq:sdot}--\ref{eq:sstar})}
  \begin{tabular}{clll}
    \toprule
    Simbolo & Significato & Unità & Nota \\
    \midrule
    $h_0$   & Modulo di incrudimento iniziale & \si{\pascal} & $h_0 \gg s$ tipicamente \\
    $a$     & Sensitività incrudimento & -- & $a > 1$ \\
    $\hat{s}$ & Resistenza di saturazione di riferimento & \si{\pascal} & \\
    $n$     & Sensitività $\sstar$ a $\epdot$ & -- & piccolo ($\sim 10^{-2}$) \\
    $s_0$   & Condizione iniziale $s(t=0)$ & \si{\pascal} & \\
    \bottomrule
  \end{tabular}
\end{table}

\subsection{Interpretazione fisica}

\begin{itemize}
  \item $s < \sstar$: incrudimento ($\sdot > 0$) --- regime di bassa temperatura
    o alta velocità di deformazione
  \item $s > \sstar$: ammorbidimento ($\sdot < 0$) --- recovery termico dominante
  \item $s = \sstar$: stato stazionario --- regime di creep secondario
\end{itemize}

% ══════════════════════════════════════════════════════════════════════════════
\section{Proprietà del Materiale \sac{}}
\label{sec:props}
% ══════════════════════════════════════════════════════════════════════════════

\subsection{Parametri Anand per \sac{} (SAC305)}

\begin{table}[H]
  \centering
  \caption{Parametri Anand per \sac{} -- valori consolidati in letteratura
    \cite{Wiese2010,Darveaux2000,JEDEC}}
  \renewcommand{\arraystretch}{1.15}
  \begin{tabular}{clSl}
    \toprule
    Simbolo & Significato & {Valore} & Unità \\
    \midrule
    $A$       & Costante pre-esponenziale & 4.0e4    & \si{\per\second} \\
    $Q/R$     & Temperatura di attivazione & 5765     & \si{\kelvin} \\
    $\xi$     & Stress multiplier & 4.0      & -- \\
    $m$       & Strain rate sensitivity & 0.303    & -- \\
    $h_0$     & Modulo di incrudimento & 3.35e9   & \si{\pascal} \\
    $\hat{s}$ & Resistenza saturazione rif. & 3.24e7   & \si{\pascal} \\
    $a$       & Sensitività incrudimento & 1.65     & -- \\
    $n$       & Sensitività $\sstar$ & 0.0074   & -- \\
    $s_0$     & Resistenza iniziale & 2.0e7    & \si{\pascal} \\
    \bottomrule
  \end{tabular}
  \label{tab:anand_params}
\end{table}

\subsection{Proprietà elastiche, termiche e meccaniche}

\begin{table}[H]
  \centering
  \caption{Proprietà fisiche del solder \sac{} (bulk, 25\,°C)}
  \renewcommand{\arraystretch}{1.15}
  \begin{tabular}{lSll}
    \toprule
    Proprietà & {Valore} & Unità & Note \\
    \midrule
    Modulo di Young $E(25\,°\text{C})$ & 50    & \si{\giga\pascal} & decresce con $T$ \\
    Coefficiente di Poisson $\nu$      & 0.35  & --       & \\
    CTE $\alpha$                       & 21    & $\mu$m/(m$\cdot$K) & isotropo \\
    Densità $\rho$                     & 7380  & \si{\kg\per\m\cubed} & \\
    Conducibilità termica $k$          & 57    & \si{\W\per\m\per\kelvin} & \\
    Calore specifico $c_p$             & 220   & \si{\J\per\kg\per\kelvin} & \\
    Temperatura solidus $T_s$          & 490   & \si{\kelvin} & 217\,°C \\
    \bottomrule
  \end{tabular}
\end{table}

\subsubsection{Dipendenza da temperatura del modulo di Young}

\begin{equation}
  E(T) = E_0 - C\cdot T
  = \SI{55e9}{\pascal} - \SI{6.5e7}{\pascal\per\kelvin}\cdot T
  \label{eq:ET}
\end{equation}

valida per $T \in [200, 490]\,\si{\kelvin}$ (al di sotto del solidus).

% ══════════════════════════════════════════════════════════════════════════════
\section{Figure TikZ}
% ══════════════════════════════════════════════════════════════════════════════

\subsection{Comportamento di \texorpdfstring{$s(t)$}{s(t)} e \texorpdfstring{$s^*$}{s*}}

\begin{figure}[H]
  \centering
  \begin{tikzpicture}
    \begin{axis}[
      width  = 0.80\textwidth,
      height = 6.5cm,
      xlabel = {Deformazione plastica cumulata $\varepsilon^{vp}_{acc}$},
      ylabel = {Resistenza $s$ [\si{\mega\pascal}]},
      xmin   = 0, xmax = 1,
      ymin   = 0, ymax = 55,
      grid   = major,
      grid style = {dashed, gray!40},
      legend pos = south east,
      tick label style = {font=\small},
      label style      = {font=\small},
    ]
      % s* (saturazione) -- orizzontale tratteggiata
      \addplot[dashed, thick, red!70!black, domain=0:1]
        {38} node[right, pos=1, font=\small] {$s^*$};

      % Curva incrudimento (s0 < s*)
      \addplot[thick, blue!70!black, domain=0:1, samples=100]
        {38 - 18*exp(-5*x)};
      \addlegendentry{$s(t)$, incrudimento ($s_0 < s^*$)};

      % Curva ammorbidimento (s0 > s*)
      \addplot[thick, orange!80!black, domain=0:1, samples=100]
        {38 + 10*exp(-6*x)};
      \addlegendentry{$s(t)$, ammorbidimento ($s_0 > s^*$)};

      % Condizioni iniziali
      \node[circle, fill=blue!70!black, inner sep=2pt] at (axis cs:0,20) {};
      \node[left, font=\small] at (axis cs:0,20) {$s_0 = \SI{20}{\mega\pascal}$};

      \node[circle, fill=orange!80!black, inner sep=2pt] at (axis cs:0,48) {};
      \node[left, font=\small] at (axis cs:0,48) {$s_0 = \SI{48}{\mega\pascal}$};
    \end{axis}
  \end{tikzpicture}
  \caption{Evoluzione della variabile interna $s$ verso la saturazione $s^*$.
    Il materiale si incrudisce se $s_0 < s^*$ e si ammorbidisce se $s_0 > s^*$.}
  \label{fig:s_evolution}
\end{figure}

\subsection{Struttura a strati del modulo SiC}

\begin{figure}[H]
  \centering
  \begin{tikzpicture}[
    layer/.style = {draw, rectangle, minimum width=7cm, minimum height=0.55cm,
                    font=\small, align=center},
    every node/.style = {font=\small}
  ]
    % Layers (bottom to top)
    \node[layer, fill=gray!25]   (hs)   at (0,0)   {Heat Sink (Al / Cu) \quad $h_{\mathrm{conv}}$};
    \node[layer, fill=yellow!20] (tim)  at (0,0.65) {TIM / thermal grease \quad $\delta\approx\SI{100}{\micro\m}$};
    \node[layer, fill=orange!15] (dbc1) at (0,1.30) {DBC bottom Cu \quad $\SI{300}{\micro\m}$};
    \node[layer, fill=blue!10]   (al2o3)at (0,1.95) {Al$_2$O$_3$ / AlN substrate \quad $\SI{380}{\micro\m}$};
    \node[layer, fill=orange!15] (dbc2) at (0,2.60) {DBC top Cu \quad $\SI{300}{\micro\m}$};
    \node[layer, fill=red!20,
          minimum height=0.80cm] (sol)  at (0,3.35) {\textbf{Solder \sac{}} \quad $\SI{100}{\micro\m}$
          \quad \textbf{[Anand model]}};
    \node[layer, fill=green!15,
          minimum height=0.90cm] (die)  at (0,4.20) {SiC die (4H-SiC) \quad $\SI{350}{\micro\m}$
          \quad $\rho(T) = \rho_0(T/T_0)^{2.3}$};
    \node[layer, fill=cyan!10]   (bond) at (0,4.90) {Bond wire / ribbon (Al/Cu)};

    % Freccia flusso di calore
    \draw[-{Stealth[length=5pt]}, thick, red!60!black]
      ($(hs.east)+(0.3,1.5)$) -- ++(0,-1.5)
      node[right, midway] {$q$ [W/m²]};

    % Freccia corrente
    \draw[-{Stealth[length=5pt]}, thick, blue!60!black]
      ($(die.east)+(0.9,0)$) -- ++(0,-4.2)
      node[right, midway, rotate=-90, anchor=south] {$I_{\mathrm{fault}}(t)$};

    % Annotazioni CTE
    \draw[decorate, decoration={brace, amplitude=4pt}]
      ($(sol.north west)+(-0.2,0)$) -- ($(sol.south west)+(-0.2,0)$)
      node[midway, left=4pt] {$\Delta\alpha$};

    \node[align=left, font=\tiny] at (-3.5,4.2) {$\alpha_{\mathrm{SiC}}=4\;\mu\mathrm{m}/(\mathrm{m\cdot K})$};
    \node[align=left, font=\tiny] at (-3.5,3.35) {$\alpha_{\mathrm{SAC}}=21\;\mu\mathrm{m}/(\mathrm{m\cdot K})$};

  \end{tikzpicture}
  \caption{Struttura a strati del modulo di potenza SiC. Il layer di solder \sac{}
    è la zona critica per l'affidabilità meccanica ed è modellato con la legge di Anand.
    La freccia a destra indica il flusso di corrente di fault durante l'apertura del SSCB.}
  \label{fig:stack}
\end{figure}

\subsection{Ciclo termico per fault ciclici (SSCB)}

\begin{figure}[H]
  \centering
  \begin{tikzpicture}
    \begin{axis}[
      width  = 0.80\textwidth,
      height = 5.5cm,
      xlabel = {Tempo [\si{\second}]},
      ylabel = {Temperatura $T_j$ [\si{\celsius}]},
      xmin   = 0, xmax = 4,
      ymin   = 20, ymax = 170,
      grid   = major,
      grid style = {dashed, gray!40},
      xtick  = {0, 0.5, 1, 1.5, 2, 2.5, 3, 3.5, 4},
      xticklabels = {0, , 1, , 2, , 3, , 4},
      tick label style = {font=\small},
      label style      = {font=\small},
    ]
      % Ciclo termico: fault ogni 1 s, picco a 150 °C, recupero a 65 °C
      \addplot[thick, blue!70!black, domain=0:4, samples=400] {
        65
        + 85*exp(-10*mod(x,1))*( (mod(x,1) < 0.5) )
        + 85*exp(-10*mod(x+0.5,1))*( (mod(x+0.5,1) < 0.5) )*0
      };

      % Approssimazione manuale con coordinate
      \addplot[thick, blue!70!black] coordinates {
        (0.00, 65) (0.02, 150) (0.10, 120) (0.20,  95) (0.40,  75) (0.60, 70)
        (1.00, 65) (1.02, 150) (1.10, 120) (1.20,  95) (1.40,  75) (1.60, 70)
        (2.00, 65) (2.02, 150) (2.10, 120) (2.20,  95) (2.40,  75) (2.60, 70)
        (3.00, 65) (3.02, 150) (3.10, 120) (3.20,  95) (3.40,  75) (3.60, 70)
        (4.00, 65)
      };

      % Linee di riferimento
      \addplot[dashed, red!60!black, domain=0:4] {150}
        node[right, pos=0.02, font=\tiny] {$T_{j,\max}=\SI{150}{\celsius}$};
      \addplot[dashed, gray,           domain=0:4] {65}
        node[right, pos=0.02, font=\tiny] {$T_{\mathrm{cool}}=\SI{65}{\celsius}$};

      % Annotazione ΔT
      \draw[{Stealth[length=4pt]}-{Stealth[length=4pt]}, thick]
        (axis cs:0.5, 65) -- (axis cs:0.5, 150)
        node[midway, right, font=\small] {$\Delta T = \SI{85}{\kelvin}$};

    \end{axis}
  \end{tikzpicture}
  \caption{Profilo termico di giunzione $T_j(t)$ durante fault ciclici del SSCB
    (frequenza \SI{1}{\hertz}, $\Delta T = \SI{85}{\kelvin}$).
    Ogni impulso corrisponde a un evento di interruzione con corrente di fault
    $I_f \approx 2\text{--}3\times I_{\mathrm{nom}}$.}
  \label{fig:thermal_cycle}
\end{figure}

\subsection{Curva \texorpdfstring{$\dot{\varepsilon}^p$}{ε̇p} vs.\ \texorpdfstring{$\sigma/s$}{σ/s}}

\begin{figure}[H]
  \centering
  \begin{tikzpicture}
    \begin{axis}[
      width  = 0.75\textwidth,
      height = 6cm,
      xlabel = {$\sigma / s$ (normalized stress)},
      ylabel = {$\dot{\varepsilon}^p / A\,e^{-Q/RT}$ (normalized strain rate)},
      xmin   = 0, xmax = 2.5,
      ymin   = 0, ymax = 12,
      restrict y to domain = 0:12,
      grid   = major,
      grid style = {dashed, gray!40},
      legend pos = north west,
      tick label style = {font=\small},
      label style      = {font=\small},
    ]
      % xi = 4, m = 0.303  -- clip to ymax
      \addplot[thick, blue!70!black, domain=0.01:1.5, samples=150]
        { (sinh(4*x))^(1/0.303) };
      \addlegendentry{$[\sinh(\xi\,\sigma/s)]^{1/m}$, $\xi=4$, $m=0.303$};

      % Regime Norton (approssimazione bassa tensione)
      \addplot[dashed, red!60!black, domain=0.01:0.6, samples=50]
        { (4*x)^(1/0.303) };
      \addlegendentry{Regime Norton: $(\xi\sigma/s)^{1/m}$};

      % Regime esponenziale (alta tensione) -- fortemente scalato per visualizzazione
      \addplot[dotted, green!50!black, thick, domain=0.8:2.5, samples=80]
        { exp(4*x/0.303) / 1e4 };
      \addlegendentry{Regime espon.: $e^{\xi\sigma/(ms)}$ (scalato $\times10^{-4}$)};

      % Linea verticale transizione
      \draw[gray, dashed] (axis cs:0.8, 0) -- (axis cs:0.8, 12)
        node[above, font=\tiny] {transizione};

    \end{axis}
  \end{tikzpicture}
  \caption{Legge di flusso di Anand normalizzata. Il termine $\sinh$ copre
    uniformemente i regimi di Norton (basse tensioni) e esponenziale (alte tensioni).}
  \label{fig:flow_law}
\end{figure}

% ══════════════════════════════════════════════════════════════════════════════
\section{Implementazione in \comsol{}}
% ══════════════════════════════════════════════════════════════════════════════

\subsection{Percorso nel modello \comsol{}}

\begin{lstlisting}[language={}, caption={Configurazione materiale Anand in COMSOL (v6.x)}]
Materials > SAC305 Solder
  Mechanical
    Elasticity: Linear Elastic
      E:   55e9 - 6.5e7*T   [Pa]     % Eq. (8)
      nu:  0.35
    Inelastic Strain: Anand Viscoplastic
      A:     4e4    [1/s]
      Q_R:   5765   [K]               % Q/R
      xi:    4.0
      m:     0.303
      h0:    3.35e9 [Pa]
      shat:  3.24e7 [Pa]
      a:     1.65
      n:     0.0074
      s0:    2.0e7  [Pa]
  Thermal Expansion
      alpha: 21e-6  [1/K]
  Heat Transfer
      k:     57     [W/(m*K)]
      Cp:    220    [J/(kg*K)]
      rho:   7380   [kg/m^3]
\end{lstlisting}

\subsection{Setup dello studio accoppiato}

\begin{lstlisting}[language={}, caption={Configurazione studio Time Dependent --- multifisica segregata}]
Study: Time Dependent
  Solver: BDF (Backward Differentiation Formula)
    Relative tolerance: 1e-4
    Max order: 2 (stiff problems)
  Nonlinear solver: Newton
    Max iterations: 25
    Damping: automatic (constant 1.0 -> backtracking)
  Segregated steps:
    Step 1: Electric Currents
      Variables: V (electric potential)
    Step 2: Heat Transfer
      Variables: T (temperature)
      Source:    Joule heating = J.E [W/m^3]
    Step 3: Solid Mechanics (Anand)
      Variables: u (displacement)
      Load:      thermal strain from step 2
\end{lstlisting}

\subsection{Variabili di post-processing chiave}

\begin{table}[H]
  \centering
  \caption{Variabili COMSOL per analisi di affidabilità del solder}
  \renewcommand{\arraystretch}{1.2}
  \begin{tabular}{lll}
    \toprule
    Variabile COMSOL & Descrizione & Utilizzo \\
    \midrule
    \texttt{solid.epe}   & Equivalent plastic strain acc.\ $\varepsilon^{vp}_{acc}$ & Coffin-Manson \\
    \texttt{solid.Ws}    & Inelastic strain energy density $\Delta W$ & Darveaux \\
    \texttt{solid.mises} & Tensione equivalente $\sigma_{vM}$ & verifica yield \\
    \texttt{solid.s}     & Variabile interna $s$ & debug modello \\
    \texttt{T}           & Temperatura & accoppiamento \\
    \texttt{ht.Qh}       & Joule heat source & sorgente termica \\
    \bottomrule
  \end{tabular}
\end{table}

% ══════════════════════════════════════════════════════════════════════════════
\section{Criteri di Vita a Fatica}
% ══════════════════════════════════════════════════════════════════════════════

\subsection{Criterio di Coffin--Manson modificato (Engelmaier)}

\begin{equation}
  N_f = \frac{1}{2}
  \left(\frac{\Delta\varepsilon^{vp}_{acc}}{2\varepsilon'_f}\right)^{1/c}
  \label{eq:coffin}
\end{equation}

con:
\begin{align}
  c &= -0.442 - 6\times10^{-4}\,T_{\mathrm{mean}}[\si{\celsius}]
       + 1.74\times10^{-2}\,\ln(1 + f)
  \label{eq:c_coeff} \\
  \varepsilon'_f &= 0.325 \quad \text{(ductility coefficient per \sac{})}
\end{align}

\subsection{Criterio energetico di Darveaux}

La densità di energia anelastica per ciclo (integrata su un ciclo stabilizzato):
\begin{equation}
  \Delta W = \oint \sigma\,\dd\varepsilon^{vp} \qquad [\si{\joule\per\m\cubed\per\text{ciclo}}]
  \label{eq:darveaux_W}
\end{equation}

La vita a fatica e la velocità di propagazione della cricca:
\begin{align}
  N_f        &= K_1 \cdot \Delta W^{K_2} \label{eq:Nf_darveaux} \\
  \frac{da}{dN} &= K_3 \cdot \Delta W^{K_4} \label{eq:crack_darveaux}
\end{align}

\begin{table}[H]
  \centering
  \caption{Costanti di Darveaux per \sac{} \cite{Darveaux2000}}
  \begin{tabular}{cSl}
    \toprule
    Costante & {Valore} & Unità \\
    \midrule
    $K_1$ & 22.4    & cicli/(J/m³)$^{K_2}$ \\
    $K_2$ & -1.02   & -- \\
    $K_3$ & 2.28e-7 & m/ciclo/(J/m³)$^{K_4}$ \\
    $K_4$ & 1.02    & -- \\
    \bottomrule
  \end{tabular}
\end{table}

% ══════════════════════════════════════════════════════════════════════════════
\section{Conclusioni}
% ══════════════════════════════════════════════════════════════════════════════

Il modello di Anand fornisce una descrizione fisicamente consistente del comportamento
viscoplastico del solder \sac{} nei moduli SiC per SSCB a media tensione.
I punti chiave sono:

\begin{enumerate}
  \item La formulazione unificata (Eq.~\ref{eq:flow}--\ref{eq:sstar}) copre sia
    il creep primario che secondario, eliminando la necessità di separare i meccanismi.
  \item I 9 parametri di Tabella~\ref{tab:anand_params} sono disponibili in letteratura
    per \sac{} e possono essere calibrati su dati di creep sperimentali tramite
    regressione non lineare (es.\ algoritmo di Levenberg-Marquardt in MATLAB).
  \item L'implementazione in \comsol{} (v5.4+) è \emph{built-in} nel modulo
    Nonlinear Structural Materials; il coupling con Joule Heating e Heat Transfer
    è gestito nativamente dal framework multiphysics.
  \item I criteri di Coffin-Manson (Eq.~\ref{eq:coffin}) e Darveaux
    (Eq.~\ref{eq:Nf_darveaux}) permettono di stimare la vita a fatica del giunto
    a partire dai risultati FEM.
\end{enumerate}

% ══════════════════════════════════════════════════════════════════════════════
\begin{thebibliography}{9}

\bibitem{Anand1985}
  L.~Anand,
  ``Constitutive equations for hot-working of metals,''
  \textit{Int.\ J.\ Plasticity}, vol.~1, no.~3, pp.~213--231, 1985.

\bibitem{Wiese2010}
  S.~Wiese and K.-J.~Wolter,
  ``Microstructure and creep behaviour of eutectic SnAg and SnAgCu solders,''
  \textit{Microelectronics Reliability}, vol.~44, pp.~1923--1931, 2004.

\bibitem{Darveaux2000}
  R.~Darveaux,
  ``Effect of simulation methodology on solder joint crack growth correlation,''
  \textit{Proc.\ ECTC}, Las Vegas, NV, 2000, pp.~1048--1058.

\bibitem{JEDEC}
  JEDEC Standard JESD22-A104,
  ``Temperature Cycling,''
  JEDEC Solid State Technology Association, 2005.

\end{thebibliography}

\end{document}
